%==============================================================================
%  CHAPTER 2 -- THEORETICAL BACKGROUND
%
%  Job of this chapter: give the reader every concept needed to follow
%  Chapters 4 and 5, and nothing else.
%
%  ONE SECTION PER AREA OF THEORY that you actually use. If a piece of theory
%  never reappears later, delete it -- it is padding, and an examiner will ask
%  why it is there.
%
%  This chapter is other people's work, so it must be fully referenced.
%  Your own contribution goes in Chapter 3.
%==============================================================================
\chapter{Theoretical background}
\label{ch:theory}

%------------------------------------------------------------------------------
\section{First area of theory}
\label{sec:theory-1}
%  For example: the thermoelastic effect, structural dynamics, image formation,
%  statistical process control -- whatever your work rests on.
%
%  Define every symbol the first time it appears. Number every equation you
%  refer to later:
%
%  \begin{equation}
%    \label{eq:example}
%    \Delta T = -K_m \, T_0 \, \Delta \sigma
%  \end{equation}
%
%  and refer to it as Eq.~\eqref{eq:example}.
%  Always write units with siunitx: \SI{9.81}{\metre\per\second\squared}.

%------------------------------------------------------------------------------
\section{Second area of theory}
\label{sec:theory-2}

%------------------------------------------------------------------------------
\section{Third area of theory}
\label{sec:theory-3}
